\documentclass{article}
\usepackage[utf8]{inputenc}
\usepackage[T1]{fontenc}
\usepackage{geometry}
\usepackage{graphicx}
\usepackage{float}
\usepackage{hyperref}
\usepackage[english,indonesian]{babel}

\title{Optimasi Pemrosesan Template LaTeX Menggunakan Pendekatan Berbasis AI}
\author{Academic Support Agent}
\date{\today}

\begin{document}
\maketitle

\begin{abstract}

\end{abstract}

\chapter{PENDAHULUAN}


\chapter{TINJAUAN PUSTAKA}


% Original content (commented out):
% \documentclass{article}

\usepackage{arxiv}

\usepackage[utf8]{inputenc} % allow utf-8 input
\usepackage[english, brazilian]{babel}
\usepackage[T1]{fontenc}    % use 8-bit T1 fonts
\usepackage{hyperref}       % hyperlinks
\usepackage{url}            % simple URL typesetting
\usepackage{booktabs}       % professional-quality tables
\usepackage{amsfonts}       % blackboard math symbols
\usepackage{nicefrac}       % compact symbols for 1/2, etc.
\usepackage{microtype}      % microtypography
\usepackage{lipsum}		% Can be removed after putting your text content
\usepackage{graphicx}
\usepackage{natbib}
\usepackage{doi}
\usepackage{float}



\title{\includegraphics[width=0.2\textwidth]{img/selo_obc.png}~ 
\\[0.25cm]
Título do Artigo  \\
\textit{Title of Paper}
}

%\date{September 9, 1985}	% Here you can change the date presented in the paper title
%\date{} 					% Or removing it

\author{ \href{https://orcid.org/0000-0000-0000-0000}{\includegraphics[scale=0.06]{orcid.pdf}\hspace{1mm}
    Autor 1} \\
	Departamento\\
	Universidade\\
	Endereço \\
	\texttt{autor1@email.com} \\
	%% examples of more authors
	\And
	\href{https://orcid.org/0000-0000-0000-0000}{\includegraphics[scale=0.06]{orcid.pdf}\hspace{1mm}
	Autor 2} \\
	Departamento\\
	Universidade\\
	Endereço\\
	\texttt{autor2@email.com} \\
	%% \AND
	%% Coauthor \\
	%% Affiliation \\
	%% Address \\
	%% \texttt{email} \\
	%% \And
	%% Coauthor \\
	%% Affiliation \\
	%% Address \\
	%% \texttt{email} \\
	%% \And
	%% Coauthor \\
	%% Affiliation \\
	%% Address \\
	%% \texttt{email} \\
}

% Uncomment to remove the date
%\date{}

% Uncomment to override  the `A preprint' in the header
%\renewcommand{\headeright}{Technical Report}
\renewcommand{\undertitle}{Textos para Discussão | Observatório do Banco Central | IE-UFRJ | \textbf{Vol.1 n.1}}
\renewcommand{\shorttitle}{Textos para Discussão | Observatório do Banco Central | IE-UFRJ}

%%% Add PDF metadata to help others organize their library
%%% Once the PDF is generated, you can check the metadata with
%%% $ pdfinfo template.pdf
\hypersetup{
pdftitle={Texto para discussão do Observatório do Banco Central},
pdfsubject={},
pdfauthor={},
pdfkeywords={},
}

\begin{document}
\maketitle

\selectlanguage{brazilian}
\begin{abstract}
	\lipsum[1]
	
	\vspace{0.5cm}
	
	\textbf{Palavras-chave: palavra-chave1. palavra-chave2. palavra-chave3.}
	
\end{abstract}

\selectlanguage{english} 
\begin{abstract}
    \lipsum[1]
    
    \vspace{0.5cm}
	
	\textbf{Keywords: keyword1. keyword2. keyword3.}
    
    \vspace{0.5cm}
	
	\textbf{JEL Classification:}
    
\end{abstract}
\selectlanguage{brazilian}


\newpage

%----------------------------------
% Introdução
%----------------------------------
\section{Introdução}

\lipsum[2-5]

%----------------------------------
% Metodologia
%----------------------------------
\section{Metodologia}

\cite{kour2014fast}

\lipsum[2-6]

\begin{figure}[H]
    \centering
    \includegraphics[scale=0.3]{img/praia-vermelha-ufrj.png}
    \caption{Campus Praia Vermelha da UFRJ. Fonte: Breno Carvalho | Agência O Globo}
    \label{fig:figura1}
\end{figure}

%----------------------------------
% Resultados e Análises
%----------------------------------
\section{Resultados e Análises}

\lipsum[2-4]

\begin{table}[H]
\centering
\caption{Exemplo de Tabela. Elaboração própria.}
\begin{tabular}{ |p{3cm}||p{3cm}|p{3cm}|p{3cm}|  }
 \hline
 \multicolumn{4}{|c|}{Country List} \\
 \hline
 Country Name     or Area Name& ISO ALPHA 2 Code &ISO ALPHA 3 Code&ISO numeric Code\\
 \hline
 Afghanistan   & AF    &AFG&   004\\
 Aland Islands&   AX  & ALA   &248\\
 Albania &AL & ALB&  008\\
 Algeria    &DZ & DZA&  012\\
 American Samoa&   AS  & ASM&016\\
 Andorra& AD  & AND   &020\\
 Angola& AO  & AGO&024\\
 \hline
\end{tabular}
\label{tab:exemplo1}
\end{table}


%----------------------------------
% Conclusão
%----------------------------------
\section{Conclusão}

\lipsum[2-4]


\bibliographystyle{unsrtnat}
\bibliography{references}   


\end{document}


\end{document}